\section{CRDT Liftability} \label{sec:crdtliftability}

% \ap{A lot of complexity for this proof is hidden in the get function. Is there a way to characterise it in more detail?}

% \ps{Explain and illustrate the key definitions}


% \begin{definition}[Sequential Type]
% A sequential type $T = (S,s_0, \Delta)$  is a triple
%     where $S$ is the set of states of an object of type $T$, $s_0 \in S $ the initial state and $\Delta: S \times Op \to S$ is the transition function mapping a state and an operation to a successor state. 
    
%     We write $s \xrightarrow{op} s'$ when $\Delta(s, op) = s'$
% \end{definition}

% % \ps{I think you should start with the definition of liftability before characterizing it. This can be extracted from the proof below.}

% \begin{definition}[Induced binary operator]
%   Given a sequential type $T$, we define the binary operator $\oplus: S \times S \to S$ induced by the transition function $\Delta$ such that for any states $s , s', s'' \in S$, $s \oplus s' = s''$ by applying the union of the operation multi-sets that lead to states $s$ and $s'$ respectively. Formally, if $s_0 \xrightarrow{op_1 \dots  \op_m} s$ and $s_0 \xrightarrow{op_1' \dots  \op_n'} s'$. Then, $s \oplus s' = \Delta ^* (s_0, \{op_1 \dots op_m\} \uplus \{op_1'\dots op_n'\})$. $\Delta ^*$ is closure of operations, applied to the initial state $s_0$.


%   \ap{try to properly specify the confluence of the $\oplus$ operator as a property of the transition system. Relate this to intention preservation possibly. }

%   The binary operator $\oplus$ is well-defined only when the transition system is \emph{confluent}. 
%   A transition system is confluent if it satisfies the Church-Rosser theorem which states that the order in which operations are evaluated is irrelevant to the final value. 
%   This allows a transition system to reorder operations if the reordering does not cause the final value of a sequence of operations to change. \ps{define this term and discuss what does it implies.}

%   \ps{$\Delta^*$ is not defined.}
  
%   \ps{In this definition, operation $op_1'$ has to be applicable to the last state reached by $op_m$. This requires the transsition function to be total.}

%   \ps{What about when there are several sequences leading to the same state? Do we take a canonical one?}  

% \ap{The original definition of CRDT also defines a query operator in the tuple but never formally specifies its details. Do we include it here?}



% \end{definition}
 \begin{definition}[Liftablity] \label{def:liftability}
A sequential type $T$ admits a $\R$-lifting if there exists a CRDT type $T\R = (\mathbb{S}, \bot,U)$ with a binary merge operator $\sqcup$ and an observation function $\get: \mathbb{S} \to S$ that materializes the value  of a CRDT and satisfies the following:

% and initial element $\bot$, and an observation function $\get : T\R \to T$, satisfying:

% \ps{Recall that this is the definition of a state-based CRDT (i.e., a merge semi-lattice.)}

% \ap{need more explanation for L2. saying that merging the states does not loose any operations when using $\oplus$.}
\begin{enumerate}[label={L\arabic*}]
\item \label{cond:lifting1} $\get(\bot) = s_0$
\item \label{cond:lifting2} For any $s_1, s_2 \in \mathbb{S}, \get(s_1 \sqcup s_2) = \get(s_1) \oplus \get(s_2)$

% with $b$ fresh \ps{definition? One idea is that neither $a \leq b$, nor the other way around, but I see that you consider something different below.} with respect to $a$:


\item \label{cond:lifting3} $\get$ is surjective: for every $s \in S$ there exists some $s' \in \mathbb{S}$ with $\get(s') = s$.
\item \label{cond:lifting4} $\sqcup$ is commutative, associative, and idempotent.
\end{enumerate}



 \end{definition}


%  \ap{Illustrate this definition with an example.}
% \ps{One thing also that matters is the assumption about $\oplus$. In general, a data type has a transition function. Maybe you can relate the two, e.g., by saying that $\oplus$ is induced by the transition function: for any two states $s$, $s'$ such that for some operation $op$ $op(s)=s'$ then $s \oplus s' = s'$.}

Consider the example of a CRDT presented in Example \ref{ex:counterCRDT}. 
A grow only counter is a lifting from a sequential counter without a decrement operation. 

% \ap{A theorem later posits that a liftable sequential type is roll-back free if the operations in $\Delta$ follow left and right moverness?}

\begin{theorem}[Liftability Characterization]
  A sequential type $T$, with a set of states $S$ and its binary operator $\oplus$ is liftable via  $\R$ iff $(S, \oplus)$ is a commutative monoid and $\oplus$ is intention-preserving.
%   \ps{Well, this is not really $T$ but instead $(\oplus,S)$.}
\end{theorem}
This proof proceeds as follows: 
First, we prove the necessity condition i.e. Assuming that T\R ~exists, we show that T is a commutative monoid. 
Second, we prove sufficiency via a construction of a lifting for T\R.

\begin{proof}[Proof (Necessity): Liftability implies commutative monoid]$ $\newline

We verify the three axioms for $(S, \oplus)$ being a commutative monoid, namely associativity and commutativity of $\oplus$ and the existence of an identity element $s_0 \in S$. 
% \medskip
% \noindent\textbf{($\Rightarrow$) Necessity: Liftability implies commutative monoid}
% \medskip
% \noindent

% We verify the three axioms of a commutative monoid in turn.

\medskip
\noindent\emph{Associativity of $\oplus$.}
Let $s, t, u \in S$ be some arbitrary states. 
Since a valid lifting exists, by surjectivity (L3), choose states $a, b, c \in \mathbb{S}$ such that $\get(a) = s$, $\get(b) = t$, $\get(c) = u$.

% and where $a, b, c$ are mutually fresh with respect to one another (this is always achievable by considering distinct replica contributions).
% \ap{I need to characterize this freshness better. This reads very vague.}

By (L2) and associativity of $\sqcup$ (L4):
\begin{align}
\get(a \sqcup (b \sqcup c)) &= \get(a) \oplus \get(b \sqcup c) = s \oplus (t \oplus u) \label{eqn:assoc1}\\ 
\get((a \sqcup b) \sqcup c) &= \get(a \sqcup b) \oplus \get(c) = (s \oplus t) \oplus u \label{eqn:assoc2}
\end{align}

Since $\sqcup$ is associative, $a \sqcup (b \sqcup c) = (a \sqcup b) \sqcup c$. 
From equations \ref{eqn:assoc1} and \ref{eqn:assoc2}, we have: $s \oplus (t \oplus u) = (s \oplus t) \oplus u$. 
Since $s, t, u$ were arbitrary, $\oplus$ is associative.
% \ps{OK}


\medskip
\noindent\emph{Commutativity of $\oplus$.}
Let $s, t \in S$ be arbitrary states. By surjectivity (L3), choose $a, b \in \mathbb{S}$ with $\get(a) = s$ and $\get(b) = t$.

By (L2): $\get(a \sqcup b) = \get(a) \oplus \get(b) = s \oplus t$

% By commutativity of $\sqcup$ (L4): $a \sqcup b = b \sqcup a$

% Applying $\get$ and (L2) to the right hand side:
% \[
% \get(b \sqcup a) = \get(b) \oplus \get(a) = t \oplus s
% \]

Since $(a \sqcup b) = (b \sqcup a)$, we conclude:
\begin{align*}
 get(a \sqcup b) & = get(b \sqcup a) \\ 
 get(a) \oplus get(b) & = get(b) \oplus get(a) \\
 s \oplus t & = t \oplus s 
\end{align*}

Since $s, t$ were arbitrary, $\oplus$ is commutative. 


\medskip
\noindent\emph{Identity element.}
Let $t \in S$ be arbitrary and choose $a \in \mathbb{S}$ with $\get(a) = t$. 
We need a state representing $s_0$ to merge with. By (L1), $\get(\bot) = s_0$. 
$\bot$ is a state which carries no updates.
% We need $\bot$ to act as a fresh contribution with respect to $a$, which holds at initialization since $\bot$ carries no updates.

So, by (L2): $\get(a \sqcup \bot) = \get(a) \oplus \get(\bot) = t \oplus s_0$

Since $\bot$ is the identity of $\sqcup$ by the CRDT definition, so $a \sqcup \bot = a$. Thus, $\get(a \sqcup \bot) = \get(a) = t$. 
% Therefore $t \oplus s_0 = t$. 
Since $\oplus$ must be intention preserving, $s_0$ must not contain any updates so that $t \oplus s_0 = t$. 
Which makes $s_0 \in S$ the identity element for the commutative monoid $(S, \oplus)$. 

By commutativity of $\oplus$, $s_0$ is a two-sided identity for $\oplus$.
% \ps{OK}

This completes the proof of necessity.
\end{proof}

\begin{proof}[Proof (Sufficiency): Commutative monoid implies liftability]$ $\newline
 Assume that for a type $T$, its set of states $S$ and binary operator $\oplus$ form a commutative monoid. We construct a valid $T\R$-lifting explicitly.

\noindent\emph{Construction:} We define a CRDT $T\R  = (\mathbb{S}, \bot, \Delta)$ with a binary merge operator $\sqcup$. 
A state in $\mathbb{S}$ is a mapping from a replica identifier to sequential type. 
A lifted CRDT has the same legal set of operations ($\Delta$) as its sequential counterpart.
Let $R$ be a set of replica identifiers, $\mathbb{S} :R \to S$ is a CRDT with an observation function $get(): \mathbb{S} \rightarrow S$ and the following properties:  

\begin{itemize}
\item \textbf{Initial element:} $get(\bot) := s_0$.
\item \textbf{Observation function:} $\get(s) := \bigoplus_{r \in R} s_r$, the fold of $\oplus$ over all replica states.
\item \textbf{Merge operator:} For states $s, s' \in \mathbb{S}, s \sqcup s' = get(s) \oplus get(s')$.
\end{itemize}

While replicas store a shared state, a replica $r$ updates its own coordinate $s_r$ in its local state. For an operation $op \in \Delta$, a state update at replica $r$ is evaluated as follows: 
% Each replica $r$ maintains only its own coordinate, so replica $r$'s update  $t$ on a state $s$ is:

\[s' = s[r \mapsto \Delta(s_r, op)]\]
% \[
% s' := \forall r' \in R, \begin{cases}
% \Delta(s_{r'}, op) &\text{if } r' = r \\
%  s_r & \text{otherwise}
% \end{cases}
% \]

% \ps{OK.}

With such a lifting, we prove conditions \ref{cond:lifting1}, \ref{cond:lifting2}, \ref{cond:lifting3} and \ref{cond:lifting4}.  

\medskip
\noindent\emph{\textbf{Condition \ref{cond:lifting1}}.} holds by the definition of our CRDT construction. Since, $\get(\bot) = s_0$. 


\medskip
\noindent\emph{\textbf{Condition \ref{cond:lifting2}.}}
Let $f, g \in T\R$ be two states of a lifted CRDT. A merge $f \sqcup g$ is defined as a fold over all replica states. 
\[
\get(f \sqcup g) = \bigoplus_{r \in R} (f_r \sqcup g_r)
\]

This pointwise fold can be expanded over all replica states as follows:

\begin{align*}
\bigoplus_{r \in R} (f_r \sqcup g_r) &= (f_1 \sqcup g_1) \oplus \dots \oplus (f_r \sqcup g_r) \\
&=   (get(f_1) \oplus get(g_1)) \oplus \dots \oplus (get(f_r) \oplus get(g_r))
\end{align*}

Individual replica contributions $f_1 \dots f_r \text{ and } g_1 \dots g_r$ are disjoint and $\oplus$ is associative, commutative and intention-preserving. So, a fold can be rewritten.
\begin{align*}
   & (get(f_1) \oplus get(g_1)) \oplus \dots \oplus (get(f_r) \oplus get(g_r)) \\
     &=   (\get(f_1) \oplus \dots \oplus get(f_r)) \oplus (\get(g_1) \oplus \dots \oplus get(g_r)) \\
    &= (\bigoplus_{r \in R} \get(f_r)) \oplus (\bigoplus_{r \in R} \get(g_r)) \\
    &= \get(f) \oplus \get(g)
\end{align*}

Thus, $\get(f \sqcup g) = get(f) \oplus \get(g)$





% fresh with respect to one another, meaning their nonzero coordinates are disjoint: for each $r$, at most one of $f(r)$, $g(r)$ is non-initial \ps{I see.}. Then:
% \[
% \get(f \sqcup g) = \bigoplus_{r \in R} (f \sqcup g)(r) = \bigoplus_{r \in R} (f(r) \oplus g(r))
% \]

% Since $\oplus$ is commutative and associative, and the supports of $f$ and $g$ are disjoint:
% \[
% = \left(\bigoplus_{r \in R} f(r)\right) \oplus \left(\bigoplus_{r \in R} g(r)\right) = \get(f) \oplus \get(g)
% \]


% \medskip
% \noindent\emph{\textbf{Condition \ref{cond:lifting3}.}}
% For any $s \in S$, fix a replica $r_0 \in R$ and define a function $f := \lambda r.\, \text{if } r = r_0 \text{ then } s \text{ else } $. Then $\get(f) = t \oplus e \oplus \cdots \oplus e = t$. So $\get$ is surjective. \ps{OK}


\medskip
\noindent\emph{\textbf{Condition \ref{cond:lifting4}}: Verification of CRDT axioms.}

\noindent\textit{Associativity of $\sqcup$} follows from associativity of $\oplus$ applied pointwise.

\noindent\textit{Commutativity of $\sqcup$} follows from the commutativity of $\oplus$.

% \ps{We might need also that $e$ is neutral; no? This is easy: $r$, $(f \sqcup \bot)(r) = f(r) \oplus \bot(r) = f(r) \oplus e = f(r)$} \ap{The proof of e (now called $s_0$) being neutral comes from the other direction where we define the CRDT type and prove that it produces a commutative monoid.}


\noindent\textit{Idempotence of $\sqcup$:} A merge via $\sqcup$ is evaluated pointwise over all replica contributions in a state. Assume states, $s, s', s'' \in \mathbb{S}\text{ such that } s\sqcup s' = s''$. 

By definition of a lifted CRDT, $s'' = get(s) \oplus get (s')$. 
When $\oplus$ is intention-preserving and $\Delta(s_0,\sigma)=s$ and $\Delta(s_0,\sigma')=s'$, $s''$ is achieved by a sequence of operations $\lambda$ such that $\Delta(s_0,\lambda)=s \oplus s'$ and $ops(\lambda)=ops(\sigma) \cup ops(\sigma')$. 
The idempotence of this sequence of operations follows from the idempotence of set union. 

This completes the construction and verification. Hence $(T, \oplus, e)$ admits a $\R$-lifting.
\end{proof}

\noindent\textbf{Conclusion}

$(T, \oplus, e)$ admits a $\R$-lifting if and only if $(T, \oplus, e)$ is a commutative monoid.
\hfill$\blacksquare$

%% \ap{present some existing CRDTs and show how this theorem applies to sequential types to get a CRDT.}
%% \ap{Further extends to the fact that the product of two comm monoids is a monoid. }

% \begin{corollary}[Non-liftability of string concatenation]
% The type $(\Sigma^*, \cdot, \varepsilon)$ admits no $\R$-lifting, since string concatenation is not commutative: $\text{``AB''} \cdot \text{``CD''} = \text{``ABCD''} \neq \text{``CDAB''} = \text{``CD''} \cdot \text{``AB''}$. By the necessity direction of the theorem, no valid $\get$ can exist. \qed
% \end{corollary}

\medskip
The result above provides some interesting observations about CRDTs.

\paragraph{Observations}
First, the fact that $(\mathbb{N},\mathit{max})$ is a commutative monoid implies the existence of a grow-only CRDT counter (see~\refsec{background}).
%
For a CRDT sequence, its states are all the partial orders $\mathcal{P}$ such that for any two partial orders $P_1=(V_0, \le_0)$ and $P_2=(V_1, \le_1)$, if (i) if $y \le_0 x$ and $x \in V_1$, then $y \le_1 x$, and (ii) the converse holds for $P_2$ wrt. $P_1$.
In that case, $(\mathcal{P}, \cup)$ where $\cup$ is the standard union operator forms a commutative monoid.

A second observation is that the Cartesian product of two commutative monoids is also a commutative monoid.
Thus, for any object $T$ with a set of states $S$ and mutators $\Delta_m$ if $(S,\Delta_m)$ can be split into a set $(S,\Delta_1), \ldots, (S,\Delta_m)$ such that each $(S,\Delta_i)$ is liftable in the sense of \reftheo{lift}, then $T$ admits a CRDT implemenation.
%
This happens with PN-counter (\refsec{background});
in this case $m=2$, $\Delta_1$ encodes increments, and $\Delta_2$ encodes decrements.

Going back to the CRDT sequence, removal is implemented by taking the Cartesian product of $(\mathcal{P}, \cup)$ together with a commutative monoid $(\mathcal{M}, \land)$ where every $m \in \mathcal{M}$ is a map from the sequence domain to $\{0,1\}$ such that $(m \land m')[x]$ equals $0$ if either $m[x]=0$ or $m'[x]=0$, and $1$ otherwise.
This decomposition correspond to the notion of tombstone, as found e.g., in WOOT~\cite{woot}.

% This is not intention preserving!
%% A third observation is that there exists a universal approach to lift sequential objects in certain cases.
%% Let $M \subseteq O$ be the mutators defined by the data type and $<$ be the common partial order over sequences in the free monoid $M^*$.
%% Suppose some linearization $\ll$ of $<$.
%% Then, denoting $\sqcup_{\ll}$ the lower upper bound operator induced by $\ll$, $(M^*,\sqcup_{\ll})$ is a commutative monoid.

\ps{It seems that there is no universal lifting, but how to prove this?}

\ps{We can establish that LWW register is not a lifting of a register.}
